\documentclass[lettersize,journal]{IEEEtran}
\usepackage{amsmath,amsfonts,amssymb}
\usepackage{algorithmic}
\usepackage{algorithm}
\usepackage{array}
\usepackage[caption=false,font=normalsize,labelfont=sf,textfont=sf]{subfig}
\usepackage{textcomp}
\usepackage{stfloats}
\usepackage{url}
\usepackage{verbatim}
\usepackage{graphicx}
\usepackage{cite}
\usepackage{hyperref}
\usepackage{orcidlink} % Để chèn icon ORCID xanh lá

\hyphenation{op-tical net-works semi-conduc-tor IEEE-Xplore}

\begin{document}

\title{Amortized Neural Estimation of Transfer Entropy for Short Time Series via Unsupervised Domain Adaptation}

% Thêm tên tác giả và ORCID
\author{Nam-Hai Nguyen-Dao~\orcidlink{0009-0003-4140-3367} and Ngoc-Son Nguyen
\thanks{N.-H. Nguyen-Dao and N.-S. Nguyen are with the School of Electrical and Electronic Engineering, Hanoi University of Science and Technology, Hanoi, Vietnam (e-mail: namhai23092005@gmail.com; son.nn6303@gmail.com).}}

\maketitle

\begin{abstract}
Transfer Entropy (TE) is a powerful model-free tool for discovering causal interactions in complex dynamical systems. However, its application to real-world biomedical signals is severely hindered by the fundamental trade-off between stationarity and sample size: ensuring stationarity requires extracting extremely short time windows (e.g., $N \le 200$), a regime where traditional and recent neural TE estimators suffer from prohibitively high variance. Furthermore, neural estimators trained on synthetic distributions inevitably encounter a critical domain gap when deployed on real-world noisy signals, leading to erroneous causal inferences. To overcome these limitations, we propose an Amortized Neural TE Estimator designed explicitly for the small-sample regime. By employing a shared network architecture that maps both source and target historical trajectories into a common representation space, we theoretically and empirically demonstrate a drastic reduction in estimation variance. Crucially, to bridge the reality gap without requiring ground-truth causal labels, we introduce an Unsupervised Domain Adaptation (UDA) strategy based on the Donsker-Varadhan bound. Extensive evaluations on both non-linear synthetic models and the real-world Fantasia cardiorespiratory dataset reveal that our method outperforms state-of-the-art baselines on short sequences ($N=10-200$) and achieves consistent directionality recovery on real signals. Our findings establish a robust, fast, and scalable framework for evaluating directed information flow in short, non-stationary physiological time series.
\end{abstract}

\begin{IEEEkeywords}
Transfer Entropy, Causality, Amortized Inference, Unsupervised Domain Adaptation, Short Time Series, Biomedical Signals.
\end{IEEEkeywords}

\section{Introduction}
% Bạn sẽ dán nội dung Introduction từ draft vào đây
\IEEEPARstart{T}{he} ability to infer directed, causal interactions from time-series data is a cornerstone...

\section{Related Work}
\subsection{Transfer Entropy and Mutual Information}
% Viết về định nghĩa TE
\subsection{Neural Estimators for Causality}
% Viết về MINE, TREET, TENDE
\subsection{Challenges in the Small-Sample Regime}
% Viết về vấn đề Variance và Domain Gap

\section{Methodology}
\subsection{Problem Formulation}
% Công thức Y_t, X_lag, Y_lag
\subsection{Amortized TE Estimator Architecture}
% Mô tả Shared Network
\subsection{Variance Reduction via Shared Networks}
% Proposition 1 & 2
\subsection{Unsupervised Domain Adaptation (UDA)}
% Quy trình fine-tune không nhãn

\section{Experiments}
\subsection{Datasets and Data Partitioning}
% Linear VAR, Periodic Coupling, Fantasia
\subsection{Training and Adaptation Setup}
% Hyperparameters, Optimizer
\subsection{Evaluation Setup}
% Metrics: MSE, Variance, Direction Accuracy, P-value

\section{Results}
\subsection{Variance Reduction on Synthetic Short Time Series}
% So sánh với KSG, Linear
\subsection{Ablation Study: The Role of Shared Network}
% So sánh với Standard MINE
\subsection{Real-World Validation: Domain Adaptation on Fantasia}
% Kết quả trước và sau khi UDA
\subsection{Statistical Significance}
% Paired Bootstrap, Wilcoxon tests

\section{Discussion}
\subsection{Mechanism of Variance Reduction}
\subsection{Robustness of UDA vs. Overfitting}
\subsection{Limitations and Future Work}

\section{Conclusion}
% Kết luận lại bài báo

\section*{Acknowledgment}
This work was supported by... (Điền quỹ tài trợ nếu có).

\bibliographystyle{IEEEtran}
\bibliography{References}

\end{document}
